% Irodov, Part II Section 2.2, original numbered problems 36--45,
% printed pp.66--67 (PDF pp.75--76). Fig. 13 in original Answers
% gives qualitative P-T and V-T diagrams, here reproduced in TikZ.

\thermoenergyproblem{36}{66}
一根两端封闭的水平圆筒被一个可无摩擦移动的
活塞分成体积均为 $V_0$ 的左右两腔。
两侧各装有压强 $p_0$、温度相同的理想气体，
初始处于平衡。
现缓慢移动活塞，使其中一腔体积变成另一腔的
$\eta$ 倍，整个过程中两腔气体\emph{等温}。
求外界移动活塞需做的功。
\exsolution
整个圆筒体积保持
$V_1+V_2=2V_0$。
题给最终体积比 $V_1/V_2=\eta$，
因此
\[
V_1=\frac{2\eta}{1+\eta}V_0,\qquad
V_2=\frac2{1+\eta}V_0.
\]
初始时两腔各自的理想气体摩尔数满足
$\nu RT=p_0V_0$。
两腔等温膨胀功之和为
\[
W_{\rm gas}=p_0V_0\ln\frac{V_1}{V_0}
+p_0V_0\ln\frac{V_2}{V_0}
=p_0V_0\ln\frac{4\eta}{(1+\eta)^2}.
\]
因 $(1+\eta)^2\ge4\eta$，
$W_{\rm gas}\le0$；
外力的准静态做功为其相反数：
\[
\boxed{A_{\rm ext}=p_0V_0
\ln\frac{(1+\eta)^2}{4\eta}}.
\]
\excheck 对 $\eta=1$，活塞不发生净位移且
$A_{\rm ext}=0$；对偏离1的比例，外力功非负。

\thermoenergyproblem{37}{67}
三摩尔理想气体初温
$T_0=273\,\mathrm K$。
先等温膨胀至初始体积的
$n=5.0$ 倍，再等容加热，
使最终压强回到最初数值。
两过程总吸热量为
$Q=80\,\mathrm{kJ}$。
求气体的热容比
$\gamma=C_p/C_V$。
\exsolution
第一段等温，
$V_1=nV_0$、$T_1=T_0$，
吸热等于气体做功：
\[
Q_1=\nu RT_0\ln n,\qquad \nu=3.
\]
第一段结束时压强为 $p_0/n$。
第二段等容使压强回到 $p_0$，
故末温为 $T_2=nT_0$。
其吸热
\[
Q_2=\nu C_V(T_2-T_1)
=\nu\frac R{\gamma-1}(n-1)T_0.
\]
总吸热为
$Q=\nu RT_0[\ln n+(n-1)/(\gamma-1)]$，
故
\[
\boxed{\gamma
=1+\frac{n-1}{Q/(\nu RT_0)-\ln n}}.
\]
代入
$\nu=3,T_0=273\,\mathrm K,n=5,Q=80\,\mathrm{kJ}$：
\[
\boxed{\gamma\simeq1.4}.
\]
\excheck 初段为等温，\emph{不}改变气体内能；
最终压力回到初压的第二段是等容加热，温度必须升至
$nT_0$。

\thermoenergyproblem{38}{67}
定性画出理想气体的
等容、等压、等温、绝热四种准静态过程曲线，
分别采用坐标
(a) $(T,p)$；
(b) $(T,V)$，其中 $T$ 为横轴。
\exsolution
从理想气体状态方程
$pV=\nu RT$ 和可逆绝热方程
$pV^\gamma=\text{常数}$ 出发：
\[
\begin{array}{c|c|c}
\text{过程}&p\text{ 对 }T&V\text{ 对 }T\\ \hline
\text{等容}&p\propto T&V=\text{常数}\\
\text{等压}&p=\text{常数}&V\propto T\\
\text{等温}&T=\text{常数}&T=\text{常数}\\
\text{绝热}&p\propto T^{\gamma/(\gamma-1)}
&V\propto T^{-1/(\gamma-1)}
\end{array}
\]
因此在 $(T,p)$ 坐标中，
等压是水平线、等容是过原点的直线、
等温是竖线、绝热为向上弯曲曲线；
在 $(T,V)$ 坐标中，
等容是水平线、等压是过原点的直线、
等温是竖线、绝热为单调下降的曲线。
以下给出两组示意图，均不表示精确比例：
\begin{center}
\begin{tikzpicture}[scale=0.78,>=Stealth]
\draw[->] (0,0)--(4.3,0) node[right] {$T$};
\draw[->] (0,0)--(0,3.2) node[above] {$p$};
\draw[thick] (0,1.2)--(4.1,1.2) node[right] {\scriptsize 等压};
\draw[thick] (0,0)--(2.8,2.8) node[right] {\scriptsize 等容};
\draw[thick] (2.1,0)--(2.1,2.95) node[above] {\scriptsize 等温};
\draw[thick,domain=0.6:2.8,samples=45]
plot (\x,{0.16*(\x)^2.6});
\node at (3.55,2.55) {\scriptsize 绝热};
\begin{scope}[xshift=7cm]
\draw[->] (0,0)--(4.3,0) node[right] {$T$};
\draw[->] (0,0)--(0,3.2) node[above] {$V$};
\draw[thick] (0,1.4)--(4.0,1.4) node[right] {\scriptsize 等容};
\draw[thick] (0,0)--(3.2,2.65) node[right] {\scriptsize 等压};
\draw[thick] (2.2,0)--(2.2,2.95) node[above] {\scriptsize 等温};
\draw[thick,domain=0.85:3.6,samples=45]
plot (\x,{2.6/(\x)^1.35});
\node at (0.9,2.9) {\scriptsize 绝热};
\end{scope}
\end{tikzpicture}
\end{center}
\excheck 绝热曲线的指数由 $\gamma>1$ 决定；
图示的坐标轴不采用一般教材常见的 $(V,p)$，
而是原题明确要求的 $(T,p)$ 和 $(T,V)$。

\thermoenergyproblem{39}{67}
一摩尔氧气初始温度
$T_0=290\,\mathrm K$，
进行可逆绝热压缩，
压强增加 $\eta=10.0$ 倍。
求：(a) 压缩后的气体温度；
(b) 外界对气体做的功。
\exsolution
取常温氧气为理想双原子气体，
绝热指数 $\gamma\simeq7/5$，
定容摩尔热容
$C_V=R/(\gamma-1)$。
可逆绝热关系
$TV^{\gamma-1}=\text{常数}$、
$pV^\gamma=\text{常数}$，
消去体积：
\[
\frac T{T_0}=
\left(\frac p{p_0}\right)^{(\gamma-1)/\gamma}.
\]
因此
\[
\boxed{T=T_0\eta^{(\gamma-1)/\gamma}
=290\cdot10^{2/7}
\simeq0.56\,\mathrm{kK}}.
\]
绝热条件 $Q=0$，
外界对气体做功
$A_{\rm on}=\Delta U=C_V(T-T_0)$，
于是
\[
\boxed{A_{\rm on}
=\frac{RT_0}{\gamma-1}
\left[\eta^{(\gamma-1)/\gamma}-1\right]
\simeq5.6\,\mathrm{kJ}}.
\]
\excheck 压缩升温且外力做正功；
气体\emph{自身}做功为负 $-A_{\rm on}$，
不能与外力做功混同。

\thermoenergyproblem{40}{67}
某质量的氮气从相同初态开始，
分别以可逆绝热方式和等温方式
把气体体积压缩 $\eta=5.0$ 倍。
求两种压缩中外力所做功的比值
$A_{\rm ad}/A_{\rm iso}$。
\exsolution
设初温 $T_0$、摩尔数 $\nu$、
绝热指数 $\gamma$。
对于绝热压缩，
末温
$T_{\rm ad}=T_0\eta^{\gamma-1}$，
外力做功等于内能增加：
\[
A_{\rm ad}=
\nu C_V(T_{\rm ad}-T_0)
=\frac{\nu RT_0}{\gamma-1}
(\eta^{\gamma-1}-1).
\]
等温压缩外力做功为
\[
A_{\rm iso}=\nu RT_0\ln\eta.
\]
相除：
\[
\boxed{\frac{A_{\rm ad}}{A_{\rm iso}}
=\frac{\eta^{\gamma-1}-1}
{(\gamma-1)\ln\eta}}.
\]
氮气 $\gamma\simeq1.40$，
$\eta=5$，得
\[
\boxed{\frac{A_{\rm ad}}{A_{\rm iso}}
\simeq1.4}.
\]
\excheck 绝热压缩过程中气体温度升高，
因此需要的外力功比等温压缩大。

\thermoenergyproblem{41}{67}
一封闭绝热圆筒内有理想气体，
中间有一个\emph{导热且可无摩擦移动}的活塞。
初始两腔体积相等，
两侧气体温度同为 $T_0$，
随后缓慢推动活塞，
使两腔体积比
$V_{\rm large}/V_{\rm small}=\eta$。
求新共同温度 $T(\eta)$，
假设各腔初始气体物质的量相同、
过程整体可逆。
\exsolution
绝热是指整体与环境没有热交换，
但外力推动活塞仍可向整个系统做功。
因活塞导热，两腔全过程温度在准静态极限下相同，
且可逆过程的整个系统熵守恒。
设两腔各有 $\nu$ 摩尔理想气体，
初始各容积为 $V_0$。
末态总体积固定 $2V_0$，且体积比为 $\eta$，
所以
\[
V_1=\frac{2\eta}{1+\eta}V_0,\quad
V_2=\frac2{1+\eta}V_0.
\]
对两腔分别使用理想气体熵变化，
\[
\Delta S_i=\nu C_V\ln\frac T{T_0}
+\nu R\ln\frac{V_i}{V_0}.
\]
总熵变化为零：
\[
2\nu C_V\ln\frac T{T_0}
+\nu R\ln\frac{V_1V_2}{V_0^2}=0.
\]
因 $R/C_V=\gamma-1$，
\[
\ln\frac T{T_0}
=\frac{\gamma-1}{2}
\ln\frac{(1+\eta)^2}{4\eta}.
\]
因此
\[
\boxed{T=T_0
\left[\frac{(1+\eta)^2}{4\eta}\right]^{(\gamma-1)/2}}.
\]
\excheck 仅靠“总体绝热”不能推出 $U$ 守恒，
因为外界移动活塞\emph{对整个系统做功}；
使用可逆等熵条件才得到本题结果。

\thermoenergyproblem{42}{67}
充有氦气的绝热容器温度为
$T=1000\,\mathrm K$，
气体通过小孔进入真空。
若容器内流速远小于出口速率，
忽略黏性，求出流速率 $v$。
\exsolution
对准稳态绝热无黏流动，
单位质量的总焓和动能守恒：
\[
h_0+\frac12v_0^2=h+\frac12v^2.
\]
题设容器内速度可忽略 $v_0\simeq0$，
并将出流气体对外静态焓的余量在理想真空出口
近似全部转为动能。
理想气体比焓为
$h=C_{p,\rm mol}T/M$，
故
\[
\frac12v^2\simeq\frac{C_{p,\rm mol}T}{M},
\quad
C_{p,\rm mol}=\frac{\gamma R}{\gamma-1}.
\]
因此
\[
\boxed{v=\sqrt{
\frac{2\gamma RT}{(\gamma-1)M}}}.
\]
氦是单原子气体，
$\gamma=5/3$、
$M\simeq4.0\times10^{-3}\,\mathrm{kg/mol}$，
\[
\boxed{v\simeq\sqrt{\frac{5RT}{M}}
\simeq3.2\,\mathrm{km/s}\approx3.3\,\mathrm{km/s}}.
\]
\excheck 这是理想绝热能量转换的流速估计，
不描述实际小孔分子逸出分布的全部动力学细节。

\thermoenergyproblem{43}{67}
一摩尔、热容比为 $\gamma$ 的理想气体
沿 $V=a/T$ 过程变化，
其中 $a$ 为正常数。
若温度上升 $\Delta T$，
求气体净吸收的热量 $Q$。
\exsolution
一摩尔理想气体有
$pV=RT$，由
$V=a/T$ 得
$p=RT/V=RT^2/a$。
沿给定过程的体积微分为
$\mathrm dV=-a\,\mathrm dT/T^2$，
所以气体做功微分
\[
\delta A=p\,\mathrm dV
=\frac{RT^2}{a}
\left(-\frac a{T^2}\mathrm dT\right)
=-R\,\mathrm dT.
\]
内能微分为
$\mathrm dU=C_V\,\mathrm dT
=R\,\mathrm dT/(\gamma-1)$。
根据热力学第一定律，
\[
\delta Q=\mathrm dU+\delta A
=R\left(\frac1{\gamma-1}-1\right)\mathrm dT.
\]
积分得
\[
\boxed{Q=\frac{R(2-\gamma)}{\gamma-1}\Delta T}.
\]
\excheck $T$ 上升使体积 $V=a/T$ 减少，
气体做功为负；不能把 $\delta A$ 当成正值。

\thermoenergyproblem{44}{67}
证明：如果理想气体所做功
与其内能的相应增量成正比，
那么它所经历的准静态过程
可写成多方方程
$pV^n=\text{常数}$。
\exsolution
设比例系数为常数 $k$，
即气体做功微分与内能变化满足
\[
\delta A=p\,\mathrm dV=k\,\mathrm dU.
\]
理想气体内能
$\mathrm dU=\nu C_V\,\mathrm dT$，
而状态方程微分
$\nu R\,\mathrm dT=p\,\mathrm dV+V\,\mathrm dp$。
用
$C_V/R=1/(\gamma-1)$：
\[
p\,\mathrm dV
=\frac{k}{\gamma-1}
(p\,\mathrm dV+V\,\mathrm dp).
\]
对于 $k\ne0$，移项得
\[
\frac{\mathrm dp}{p}
=-\left(1-\frac{\gamma-1}{k}\right)
\frac{\mathrm dV}{V}.
\]
令多方指数
$n=1-(\gamma-1)/k$，
积分为
$\ln p+n\ln V=\text{常数}$，
即
\[
\boxed{pV^n=\text{常数},\qquad
n=1-\frac{\gamma-1}{k}}.
\]
\excheck 如果 $k=0$，
则 $\delta A=0$ 对应体积不变的等容过程，
可视为多方指数 $n\to\infty$ 的极限；
证明中的除以 $k$ 步骤不适用于这一退化情形。

\thermoenergyproblem{45}{67}
一摩尔理想气体的绝热指数为 $\gamma$，
沿多方过程 $pV^n=\text{常数}$
变化。
求过程摩尔热容 $C_n$，
并判断何种多方指数 $n$ 下 $C_n<0$。
\exsolution
多方关系微分为
\[
V\,\mathrm dp+n p\,\mathrm dV=0.
\]
一摩尔理想气体状态方程
$RT=pV$，两边求微分得
\[
R\,\mathrm dT=p\,\mathrm dV+V\,\mathrm dp
=(1-n)p\,\mathrm dV.
\]
因此沿此过程
$\delta A=p\,\mathrm dV=
R\,\mathrm dT/(1-n)$，
而内能微分
$\mathrm dU=C_V\,\mathrm dT
=R\,\mathrm dT/(\gamma-1)$。
按 $C_n=\delta Q/\mathrm dT$，
\[
C_n=
\frac{R}{\gamma-1}+\frac{R}{1-n}
=\boxed{\frac{R(n-\gamma)}
{(n-1)(\gamma-1)}}.
\]
对通常 $\gamma>1$ 且 $n\ne1$，
它为负数当且仅当
\[
\boxed{1<n<\gamma}.
\]
\excheck $n=1$ 是等温过程，
$\mathrm dT=0$ 而可发生热量交换，
因而“摩尔热容 $\delta Q/\mathrm dT$”
在这一点不具通常有限定义；
$n=\gamma$ 为绝热过程，$C_n=0$。
